\section{Trust Modeling and the Trust-Aware GNN}
\label{sec:trustgnn}

This chapter describes the learning and inference components that sit on top
of the pruned communication graph produced by the QUBO layer. The design
centers on a per-node trust model that converts observed mission outcomes
into a slowly evolving reputation signal, and a lightweight graph neural
network that scores candidate rovers against task requirements using exactly
the pruned topology as its input graph. We first state the deployment
constraint that shapes every architectural choice, then describe the node
features, the trust model, the message-passing layer, the assignment rule,
and finally the empirical evidence: what the trust model tracks, what the
GNN contributes, and where our training attempt produced an honest null
result.

\subsection{Deployment constraint: mist, fog, and cloud tiers}

The target platform is a three-tier architecture. Rovers and their immediate
wireless neighbors form the mist tier; a nearby relay or base station with
Raspberry Pi class hardware forms the fog tier; and an offboard server with
unbounded compute forms the cloud tier. This split imposes a hard asymmetry
on the learning system. Training, which requires storing trajectories,
running many epochs of gradient updates, and evaluating held-out folds,
happens exclusively in the cloud tier. The fog tier, where assignment
decisions must be made under latency and power budgets, only
\emph{evaluates} the model. Every forward pass on the fog tier must fit in
the memory and instruction budget of a small single-board computer, and must
not depend on any deep learning framework.

Consequently, the entire inference path is implemented in pure NumPy with
tiny dense matrices. Node feature vectors have four components; hidden
states have the same dimensionality as the node count, which in our missions
never exceeds the low tens; and the message-passing round is a single
matrix-vector product followed by an elementwise nonlinearity. There are no
learned embeddings, no attention heads, and no multi-layer stacks. This is a
deliberate restriction: the question we ask is not how expressive a GNN we
can afford, but whether a minimal one-round architecture already captures
the neighborhood information that matters for task allocation.

\subsection{Node features}

Each rover is described by a four-dimensional feature vector, with every
component normalized to the interval $[0,1]$ so that weights in the scoring
head are directly comparable. Table~\ref{tab:nodefeatures} lists the
features.

\begin{table}[t]
  \centering
  \caption{Node features used by the trust-aware GNN. All values are
  normalized to $[0,1]$ and refreshed every mission cycle.}
  \label{tab:nodefeatures}
  \begin{tabular}{lll}
    \toprule
    Feature & Definition & Interpretation \\
    \midrule
    Trust & EWMA reputation from the trust model & Expected task reliability \\
    Battery & Current charge relative to capacity & Energy headroom \\
    Available compute & $1 - \text{load}$ & Free processing capacity \\
    Capability & Hardware class, normalized & Sensor and actuator quality \\
    \bottomrule
  \end{tabular}
\end{table}

Three of the four features are straightforward instrumentation: battery is
read from the power subsystem, available compute is one minus the current
scheduling load, and capability is a static property of the rover's hardware
class. The fourth feature, trust, is the output of the trust model described
next, and it is the component that couples inference across mission cycles.

\subsection{The trust model}

The trust model maintains a scalar reputation per node, updated once per
mission cycle from two evidence sources: the observed outcomes of tasks the
node was assigned, and link-level evidence gathered from communication with
the node. Updates follow an exponential weighted moving average, so recent
outcomes dominate the estimate while older outcomes decay geometrically. Two
additional mechanisms keep the estimate well behaved in a dynamic network.
First, \emph{trust aging} lets evidence fade over time even in the absence
of new observations, so that a rover that has been silent for many cycles
drifts toward a neutral prior rather than retaining a stale extreme rating.
Second, a \emph{relax} operation handles nodes that received no assignment
in a cycle: since no outcome can be observed for them, their trust is
nudged toward the prior instead of being left frozen, which prevents
unassigned nodes from being permanently locked out or permanently
privileged by an old rating.

The trust estimate closes the loop between cycles. It feeds the next cycle
in two places simultaneously: as the trust component of the QUBO edge
features, where it influences which communication links survive pruning, and
as the trust component of the GNN node features, where it influences which
rovers are scored highly for which tasks. A rover that performs well is
therefore both more likely to remain well connected and more likely to be
assigned work, and the reverse holds for a rover that fails. This coupling
is intentional, but it also means that errors in the trust estimate
propagate, which is why we measure its quality directly.

The ground truth in our simulator is a hidden, spatially correlated
reliability field per node; observed outcomes are sampled from it. We
measure how well the trust model tracks this hidden quantity using the
Spearman rank correlation between estimated trust and true reliability
across nodes, computed per cycle. Trust tracks reliability closely, and the
benefit is causal rather than merely correlational: enabling trust updates,
as opposed to freezing trust at its initial stale value, yields a per-cycle
staleness margin of $+13.75$ percentage points in mission success. Fresh
trust estimates are therefore not a cosmetic refinement; they are a major
contributor to end-to-end performance.

\subsection{One round of message passing on the pruned graph}

The GNN layer performs exactly one round of message passing, and it operates
\emph{only on the pruned graph}. The edge set selected by the QUBO layer is
literally the input graph of the GNN: links that were pruned away carry no
messages, regardless of how physically close the two rovers are. This
coupling is the defining feature of our pipeline. The QUBO layer decides
which information pathways exist, and the GNN layer decides what to conclude
from the information that flows along them.

Formally, let $A$ be the weighted adjacency of the pruned graph, with edge
weights derived from link quality and trust, and let $x_i \in [0,1]^4$ be
the feature vector of node $i$. The hidden state of node $i$ after one round
is
\begin{equation}
  h_i = \sigma\!\left( W_{\mathrm{self}}\, x_i
    + \sum_{j \in \mathcal{N}(i)} a_{ij}\, W_{\mathrm{agg}}\, x_j \right),
  \label{eq:messagepassing}
\end{equation}
where $\sigma$ is the elementwise sigmoid, $\mathcal{N}(i)$ is the neighbor
set in the pruned graph, $a_{ij}$ is the adjacency weight of the surviving
edge, and $W_{\mathrm{self}}, W_{\mathrm{agg}} \in \mathbb{R}^{4 \times 4}$
are configurable weight matrices. A readout layer then scores each
(node, task) pair by comparing the hidden state against the task's
requirement vector, producing a scalar suitability score per pair.
Figure~\ref{fig:messagepassing} illustrates one round on a small pruned
graph.

\begin{figure}[t]
  \centering
  \begin{tikzpicture}[
      node distance=10mm and 14mm,
      rover/.style={circle, draw, minimum size=8mm, inner sep=0pt, font=\small},
      op/.style={rectangle, rounded corners, draw, minimum width=16mm,
        minimum height=8mm, font=\small},
      score/.style={rectangle, draw, minimum width=14mm, minimum height=8mm,
        font=\small},
      edge/.style={-{Stealth[length=2mm]}, thick},
      pruned/.style={dashed, gray},
      lbl/.style={font=\footnotesize\itshape}
    ]
    % Source nodes (pruned graph neighborhood)
    \node[rover] (v0) {$v_0$};
    \node[rover, above right=of v0] (v1) {$v_1$};
    \node[rover, below right=of v0] (v2) {$v_2$};
    \node[rover, right=of v2, gray] (v3) {$v_3$};

    % Aggregation and readout
    \node[op, right=24mm of v0] (agg) {aggregate};
    \node[op, right=12mm of agg] (sig) {$\sigma(\cdot)$};
    \node[op, right=12mm of sig] (read) {readout};
    \node[score, right=12mm of read] (s) {scores};

    % Surviving edges
    \draw[edge] (v1) -- node[lbl, above, sloped] {$a_{10}\,x_1$} (agg);
    \draw[edge] (v2) -- node[lbl, below, sloped] {$a_{20}\,x_2$} (agg);
    \draw[edge] (v0) -- node[lbl, above] {$W_{\mathrm{self}}\,x_0$} (agg);
    \draw[edge] (agg) -- (sig);
    \draw[edge] (sig) -- node[lbl, above] {$h_0$} (read);
    \draw[edge] (read) -- node[lbl, above] {task req.} (s);

    % Pruned edge: v3 was cut by the QUBO layer
    \draw[pruned] (v3) -- node[lbl, above, sloped] {pruned} (agg);

    % Stage labels
    \node[lbl, above=6mm of v1] {pruned graph};
    \node[lbl, above=6mm of agg] {one round};
    \node[lbl, above=6mm of s] {(node, task)};
  \end{tikzpicture}
  \caption{One message-passing round for node $v_0$ on a pruned graph.
  Neighbors $v_1$ and $v_2$ survive QUBO pruning and contribute weighted
  features to the aggregation; the link to $v_3$ was pruned and carries no
  message. The aggregated signal passes through a sigmoid and a readout
  that scores the node against task requirements.}
  \label{fig:messagepassing}
\end{figure}

The default weight matrices are set by hand to implement the intended
allocation policy: prefer rovers with high trust, high battery, low load
(that is, high available compute), and a strong, trustworthy neighborhood.
The neighborhood term is what distinguishes the GNN from a purely per-node
scorer: a rover adjacent to reliable partners inherits some of their
credibility, which is exactly the bias we want when task success depends on
multi-hop delivery through the pruned topology.

\subsection{Assignment rule}

Scores are converted into assignments by a capacity-aware greedy matching.
Pairs are considered in decreasing order of score; a pair is accepted if the
rover still has free capacity and the task is still unassigned. The default
capacity is one task per node per cycle, which matches the hardware
assumption that a rover executes one mission task at a time. Greedy matching
is not optimal in general, but it is fast, deterministic given the scores,
and easy to audit, and our baselines use the same matching rule so that
comparisons isolate the scoring functions rather than the matcher.

\subsection{Cloud readout training: a null result}

Although the fog tier only evaluates, the readout layer can be trained
offline in the cloud tier on trajectories collected from past missions, and
the trained weights can then be shipped back to the fog tier. We implemented
this and evaluated it on held-out missions. The result is a genuine null:
the trained readout did not outperform the fixed hand-set weights. The
measured difference was approximately $-3.9$ percentage points of mission
success in favor of the fixed weights, with a paired permutation $p$-value
of $0.49$, well within noise. We report this as a null result rather than
discarding it. Our interpretation is twofold. First, with only four input
features and a one-round architecture, the hand-set policy already sits near
the ceiling of what the model class can express, so gradient training has
little room to add and ample room to overfit the training missions. Second,
the dominant signal, trust, is computed upstream by the trust model; the
readout mostly has to avoid corrupting it. We therefore ship the fixed
weights and keep cloud training as an optional, disabled path.

\subsection{Ablation: what actually carries the performance}

To attribute mission success to the right component, we ran a four-arm
ablation on 10-rover missions with paired statistics across seeds. The arms
are the full pipeline (GNN scoring on the pruned graph), a no-graph readout
that scores nodes from their own features without message passing, a
trust-only arm that scores nodes by trust alone, and random assignment.
Table~\ref{tab:ablation} reports the results.

\begin{table}[t]
  \centering
  \caption{Ablation on 10-rover missions. Mission success rates with paired
  statistics against the full GNN pipeline: bootstrap 95\% confidence
  intervals on the success-rate difference and permutation $p$-values.
  The no-graph difference is not statistically significant.}
  \label{tab:ablation}
  \begin{tabular}{lllll}
    \toprule
    Arm & Success & $\Delta$ vs.\ GNN (pp) & 95\% CI & $p$ \\
    \midrule
    GNN (full pipeline) & 39.8\% & $0$ & & \\
    No-graph readout & 38.1\% & $+1.7$ & not significant & n.s. \\
    Trust-only & 31.7\% & $+8.13$ & $[+4.2,\, +11.7]$ & $0.016$ \\
    Random & 26.0\% & $+13.8$ & & \\
    \bottomrule
  \end{tabular}
\end{table}

Three conclusions follow. First, the full pipeline beats trust-only scoring
by $+8.13$ percentage points, with a bootstrap 95\% confidence interval of
$[+4.2, +11.7]$ and a permutation $p$-value of $0.016$: the feature mix and
the gating structure deliver a real, statistically supported gain over using
trust alone. Second, the pipeline beats random assignment by $+13.8$
percentage points, confirming that the system extracts substantial signal
from the features at all. Third, and most instructive, the gap between the
GNN and the no-graph readout is only $+1.7$ percentage points and is
\emph{not} statistically significant. The single message-passing round on
the pruned graph does not measurably improve over per-node scoring at this
mission scale.

We state this plainly rather than softening it: at the scale we tested, the
heavy lifting is done by the trust model and by the gating of assignment
through the pruned topology, not by neighborhood aggregation itself. The
trust model's $+13.75$ point staleness margin and the $+8.13$ point margin
over trust-only scoring are the load-bearing results of this chapter. The
GNN layer contributes a principled way to fuse features and a natural
interface to the QUBO-pruned graph, and it may matter more at larger scales
or in missions where multi-hop delivery dominates; whether it does is an
empirical question we leave to the scaling experiments. What the present
evidence supports is a system whose intelligence is concentrated in a simple,
auditable trust estimate, with the neural component acting as a restrained
and honest scoring head on top of it.
